\Subhead{2.8.2}{Уравнение Бюргерса. Генерация гармоник. Нелинейное поглощение, насыщение и детектирование}

Слабая диссипация регуляризует крутой фронт:
\begin{equation}\label{eq:burgers}
 u_t+u u_x=\nu u_{xx},\qquad \nu>0.
\end{equation}
Это нормированная форма Бюргерса; в акустических координатах коэффициент нелинейности и длина распространения переводятся из предыдущего подвопроса, а эффективная диффузия содержит вязкость и теплопроводность. Подстановка $u=-2\nu\partial_x\ln\psi$ превращает уравнение в $\psi_t=\nu\psi_{xx}$ (преобразование Коула–Хопфа), что даёт способ построения решения.

Для стационарного бегущего фронта между $u_L>u_R$ скорость $V=(u_L+u_R)/2$ и
\begin{equation}
 u(x,t)=V-\frac{u_L-u_R}{2}\tanh\left[\frac{(u_L-u_R)(x-Vt)}{4\nu}\right].
\end{equation}
Нелинейность сжимает фронт, диссипация его расширяет; толщина порядка $\nu/(u_L-u_R)$. При убывающих полях или периодических условиях умножение уравнения на $u$ даёт
\begin{equation}
 \frac{\dd}{\dd t}\int\frac{u^2}{2}\dd x=-\nu\int u_x^2\dd x.
\end{equation}
Нелинейность непосредственно не уничтожает энергию: она переносит её к большим волновым числам, где потери быстрее. Поэтому с ростом амплитуды эффективное затухание усиливается; после формирования пилообразного профиля амплитуда ограничивается нелинейными потерями.

На малых расстояниях $u\simeq u_a\sin\omega\tau+(\beta xu_a^2\omega/2c_0^2)\sin2\omega\tau$: уже квадратичный член рождает вторую гармонику. При сильном искажении нужны многие гармоники. Детектирование — появление медленной компоненты из разностных частот; например, усреднение $u^2$ даёт $|A(\tau)|^2/2$ для модулированной несущей. Наличие такой нелинейной компоненты следует отличать от постоянного акустического течения, которому нужны диссипация и отдельный баланс импульса.
